% Quatrième de couverture : résumé et abstract (environ 10 lignes chacun, même contenu)
\clearpage
\thispagestyle{empty}
\noindent\textbf{\large Résumé}\par\medskip
\noindent Ce projet de fin d'études, réalisé au sein de \entreprise, porte sur la conception et la réalisation de \nomprojet, une plateforme SaaS multi-cliniques de prise de rendez-vous en ligne et de suivi des patients. La solution permet aux patients de réserver un créneau sans appel téléphonique, aux secrétariats de gérer les agendas des praticiens et aux cliniques de réduire les rendez-vous non honorés grâce à des rappels automatiques par SMS et par e-mail. Elle repose sur une architecture Spring Boot et React, une base PostgreSQL isolée par clinique, une authentification centralisée avec Keycloak et un déploiement conteneurisé sur Kubernetes, livré par une chaîne d'intégration continue. \todo{Compléter avec les résultats chiffrés.}

\noindent\textbf{Mots-clés :} SaaS, e-santé, prise de rendez-vous, multi-location, Spring Boot, React, PostgreSQL, Keycloak, Kubernetes, Scrum.

\vspace{1.5cm}
\noindent\textbf{\large Abstract}\par\medskip
\noindent This end-of-studies project, carried out at \entreprise, covers the design and implementation of \nomprojet, a multi-clinic SaaS platform for online appointment booking and patient follow-up. Patients book a slot without a phone call, front desks manage the practitioners' calendars, and clinics reduce missed appointments thanks to automatic SMS and e-mail reminders. The solution relies on a Spring Boot and React architecture, a PostgreSQL database isolated per clinic, centralised authentication with Keycloak and a containerised deployment on Kubernetes, delivered by a continuous integration pipeline. \todo{Add the measured results.}

\noindent\textbf{Keywords:} SaaS, e-health, appointment booking, multi-tenancy, Spring Boot, React, PostgreSQL, Keycloak, Kubernetes, Scrum.
